%% ma: L10-B04-0003
%% lop: 10
%% chuong: 2
%% bai: 4
%% dang: 10.4.2
%% loai: TLN
%% muc: VD
%% dap_an: "3"
%% nguon: "(NBV)-ĐỀ SỐ 18-MỖI NGÀY 1 ĐỀ THI 2025.pdf (D18, MỖI NGÀY 1 ĐỀ - Đề số 18), Phần III Câu 1"
%% kien_thuc: "Bài 3, Bài 4 (hệ bất phương trình bậc nhất hai ẩn, miền nghiệm, giá trị nhỏ nhất của biểu thức trên miền đa giác)"
%% trang_thai: chinh-thuc
%% ghi_chu: "Đáp án gốc 3 đúng (giải độc lập bằng scipy.linprog: nhỏ nhất 150 tại (0; 1,5)). Bỏ ảnh trang trí (cuộc thi bữa ăn dinh dưỡng). Đề không có hình toán học. Đã duyệt 19/10/2026."
\begin{ex}
Trong một cuộc thi về ``bữa ăn dinh dưỡng'', ban tổ chức yêu cầu để đảm bảo lượng dinh dưỡng hằng ngày thì mỗi gia đình có $4$ thành viên cần ít nhất $900$ đơn vị prôtêin và $400$ đơn vị lipit trong thức ăn hằng ngày. Mỗi kg thịt bò chứa $800$ đơn vị prôtêin và $200$ đơn vị lipit, $1$ kg thịt heo chứa $600$ đơn vị prôtêin và $400$ đơn vị lipit. Biết rằng người nội trợ chỉ được chi tối đa $200$ nghìn đồng để mua thịt, $1$ kg thịt bò giá $200$ nghìn đồng, $1$ kg thịt heo giá $100$ nghìn đồng. Người nội trợ nên mua $x$ kg thịt bò và $y$ kg thịt heo để chi phí thấp nhất cho khẩu phần thức ăn mà vẫn đảm bảo chất dinh dưỡng. Khi đó hãy tìm $x+2y$.
\shortans{3}
\loigiai{Chi phí là $F(x;y)=200x+100y$ (nghìn đồng). Các ràng buộc: $800x+600y\ge900$, $200x+400y\ge400$, $200x+100y\le200$, $x\ge0$, $y\ge0$, tức
\[\begin{cases}8x+6y\ge9\\x+2y\ge2\\2x+y\le2\\x\ge0,\ y\ge0.\end{cases}\]
Miền nghiệm là tứ giác có các đỉnh $D(0;1{,}5)$, $A(0{,}6;0{,}7)$, $B\left(\dfrac23;\dfrac23\right)$, $C(0;2)$. Ta có $F(D)=150$, $F(A)=190$, $F(B)=200$, $F(C)=200$. Chi phí nhỏ nhất bằng $150$ nghìn đồng tại $x=0$, $y=1{,}5$. Vậy $x+2y=3$.}
\end{ex}
